% Gesamttest «Funktionstypen und Graphen» — Quelle des PDFs (python3 scripts/build-lp-pdf.py funktionstypen).
% Musterlösung und Raster: bewertungspaket.tex. Alle Werte mit python3 nachgerechnet (scripts/lp/funktionstypen/zahlen.py).
% Vorgabe ~/TODO-BM-Vorkurs-LP03-15.md §4: Graphen und Gleichungen mit Testpunkten zuordnen (G1, G2), proportional,
% antiproportional, linear mit Anfangswert (G3, G4), Weg-Zeit-Graph lesen und zu einer Beschreibung wählen (G5, G6),
% zurückgelegt gegen verbleibend und fehlerhaftes Bewegungsdiagramm kritisieren (G7, G8). 25 Punkte, rund 30 Minuten.
\documentclass[11pt]{article}
\usepackage{../lp-druck}
\usepackage{ft-gitter}
\renewcommand{\lpname}{Funktionstypen und Graphen}
\renewcommand{\lpteilgebiet}{Funktionen ·}
\renewcommand{\lpfuss}{Gesamttest · 25 Punkte}
\begin{document}
\lpkopf{Gesamttest}{%
  \vspace{4pt}\noindent\begin{tabularx}{\linewidth}{@{}X X@{}}
  Code (kein Name): \hrulefill & Punkte: \hrulefill\ / 25
  \end{tabularx}}

\begin{lpinfo}{Anleitung}
\textbf{Zeit:} rund 30 Minuten. \textbf{Hilfsmittel:} Lineal, kein Taschenrechner.
Schreib zu jeder Aufgabe den \textbf{Rechenweg} oder eine \textbf{Begründung} auf; bewertet wird der Weg.
Danach: selbst mit dem \emph{Bewertungspaket} bewerten oder die Lösung samt diesem Blatt scannen bzw.\ fotografieren
(ohne Namen und Standort) und mit dem Bewertungspaket einer KI geben.
\end{lpinfo}

\lpteil{Teil A · Vier Graphen-Familien}{Kapitel 1 · 7 Punkte}

\begin{lpaufgabe}{G1}{Zuordnen}{4 P}
Ordne jedem Graphen A bis D eine Gleichung zu und begründe jede Zuordnung mit einem Merkmal oder einem Testpunkt.
\[ y = 2x + 1 \qquad y = 0.5x^2 \qquad y = 2\sqrt{x} \qquad y = \dfrac{4}{x} \]
\begin{center}
\kleingraph{A}{\draw[lpblau,line width=1pt,domain=0:5,samples=60] plot (\x,{2*sqrt(\x)});}\hfill
\kleingraph{B}{\draw[lpblau,line width=1pt,domain=-2:5,samples=60] plot (\x,{0.5*\x*\x});}\hfill
\kleingraph{C}{\draw[lpblau,line width=1pt,domain=0.42:5,samples=60] plot (\x,{4/\x}); \draw[lpblau,line width=1pt,domain=-2:-1.3,samples=20] plot (\x,{4/\x});}\hfill
\kleingraph{D}{\draw[lpblau,line width=1pt] (-2,-3) -- (4,9);}
\end{center}
\lpplatz{4}
\end{lpaufgabe}

\begin{lpaufgabe}{G2}{Wertetabelle}{3 P}
Fülle die Wertetabelle von $y = \dfrac{8}{x}$ aus und skizziere den Graphen für $x > 0$. Warum gibt es bei $x = 0$ keinen Wert?
\[ \begin{array}{c|cccc} x & 1 & 2 & 4 & 8 \\ \hline y & \phantom{00} & \phantom{00} & \phantom{00} & \phantom{00} \end{array} \]
\lpplatz{5}
\end{lpaufgabe}

\lpteil{Teil B · Proportional, antiproportional, linear}{Kapitel 2 · 6 Punkte}

\begin{lpaufgabe}{G3}{Zusammenhänge}{3 P}
Welcher Zusammenhang ist es: proportional, antiproportional oder linear mit Anfangswert? Begründe kurz.
\begin{enumerate}[label=(\alph*),leftmargin=*,itemsep=3pt]
\item Preis $y$ für $x$ kg Kartoffeln zu $2.40$ Fr. pro kg.
\item Fahrzeit $y$ für $120$ km bei $x$ km/h.
\item Handykosten $y$: $20$ Fr. im Monat und $0.10$ Fr. pro Minute.
\end{enumerate}
\lpplatz{4}
\end{lpaufgabe}

\begin{lpaufgabe}{G4}{Antiproportional}{3 P}
$6$ gleich schnelle Personen brauchen für eine Arbeit $10$ Tage.
\begin{enumerate}[label=(\alph*),leftmargin=*,itemsep=3pt]
\item Wie lange brauchen $4$ Personen?
\item Schreib die Gleichung für die Dauer $y$ bei $x$ Personen.
\item Wie lange brauchen $12$ Personen?
\end{enumerate}
\lpplatz{4}
\end{lpaufgabe}

\lpteil{Teil C · Weg-Zeit-Graphen}{Kapitel 3 · 6 Punkte}

\begin{lpaufgabe}{G5}{Velofahrt}{4 P}
Der Graph zeigt die Entfernung $s$ eines Velofahrers von zu Hause.
\begin{center}
\wegzeit{5}{9}{5}{200}{$s$ [m]}{\draw[lpblau,line width=1.2pt] (0,0) -- (2,3) -- (3,3) -- (5,9);}
\end{center}
\begin{enumerate}[label=(\alph*),leftmargin=*,itemsep=3pt]
\item Wie schnell fährt er in den ersten $10$ Minuten (in m/min)?
\item Was passiert von $10$ bis $15$ Minuten?
\item Wie schnell fährt er von $15$ bis $25$ Minuten?
\item In welchem Abschnitt fährt er am schnellsten? Woran siehst du das im Graphen?
\end{enumerate}
\lpplatz{5}
\end{lpaufgabe}

\begin{lpaufgabe}{G6}{Welcher Graph?}{2 P}
Tim fährt mit gleichem Tempo von zu Hause weg, macht eine Pause und fährt dann langsamer nach Hause zurück.
Welcher Graph passt? Begründe, warum die beiden anderen nicht passen.
\begin{center}
\skizze{A}{(0,0) -- (2,3.5) -- (3.5,3.5) -- (6.5,0)}\hfill
\skizze{B}{(0,0) -- (2,3.5) -- (3.5,3.5) -- (4.5,0)}\hfill
\skizze{C}{(0,0) -- (2,3.5) -- (3.5,3.5) -- (6.5,4.2)}
\end{center}
\lpplatz{4}
\end{lpaufgabe}

\lpteil{Teil D · Zurückgelegt oder noch entfernt}{Kapitel 4 · 6 Punkte}

\begin{lpaufgabe}{G7}{Wanderung}{3 P}
Eine Wanderung ist $12$ km lang. Du gehst gleichmässig $4$ km/h.
\begin{enumerate}[label=(\alph*),leftmargin=*,itemsep=3pt]
\item Schreib den zurückgelegten Weg $s$ und die Entfernung $d$ zum Ziel als Gleichung ($t$ in Stunden).
\item Wie gross sind $s$ und $d$ nach $2$ Stunden?
\item Wann bist du noch $3$ km vom Ziel entfernt?
\end{enumerate}
\lpplatz{5}
\end{lpaufgabe}

\begin{lpaufgabe}{G8}{Diagramm prüfen}{3 P}
Jana zeichnet den \textbf{zurückgelegten Weg} ihres Schulwegs. Das Diagramm hat zwei Fehler. Beschreibe beide und begründe, warum sie unmöglich sind.
Skizziere daneben ein mögliches Diagramm mit einer Pause.
\begin{center}
\wegzeit{6}{4}{5}{200}{zurückgelegt [m]}{\draw[lpblau,line width=1.2pt] (0,0) -- (1,2) -- (1,3.5) -- (2,3.5) -- (3,2.5);}
\end{center}
\lpplatz{4}
\end{lpaufgabe}
\end{document}
